% ==============================================================================
% PAGE 35: FABRICATION - PART 2 (Numbered Page 27)
% ==============================================================================
\subsubsection*{5. Power Supplying Unit (Solar Panel \& Battery)}
\noindent
\begin{spacing}{1.2}
{\fontsize{10.2}{14.5}\selectfont
A small solar panel (12V) is mounted on the base frame or beside the dish.\\[2.5mm]
The panel converts sunlight into electrical energy to charge a 6V battery.\\[2.5mm]
The battery supplies power to the motor driver, sensors, and control circuit for automatic tracking.\\[2.5mm]
A blocking diode is used to ensure one-way current flow from the panel to the battery, preventing discharge at night.
}
\end{spacing}

\vspace{4mm}
\subsubsection*{6. Water Circulation System}
\noindent
\begin{spacing}{1.2}
{\fontsize{10.2}{14.5}\selectfont
Water from the storage tank which is placed at a certain height for direct flow of water to the inlet of receiver tube.\\[2.5mm]
Heated water exits from the outlet of the receiver tube and is collected in another container.
}
\end{spacing}

\vspace{4mm}
\subsubsection*{7. Assembly \& Testing}
\noindent
\begin{spacing}{1.2}
{\fontsize{10.2}{14.5}\selectfont
All components (frame, dish, receiver, motor, sensors, battery) are assembled.\\[2.5mm]
The system is tested under sunlight to check tracking accuracy and water temperature rise. Temperature readings are taken at regular intervals for performance analysis.
}
\end{spacing}

\vspace{4mm}
\subsubsection*{8. Finishing}
\noindent
\begin{spacing}{1.2}
{\fontsize{10.2}{14.5}\selectfont
Wiring is neatly arranged using cable ties.\\[2.5mm]
Components are painted or coated to prevent corrosion.\\[2.5mm]
Labels are added for inlet, outlet, and electrical parts for clarity.
}
\end{spacing}
\clearpage
